\documentclass[10pt,a4paper]{amsart}
\usepackage{amsmath,amssymb,amsthm,graphicx}
% Retained preamble for the Stacks Project chapter snapshot in this directory.
% The chapter's own preamble.tex is a copy of the shared Stacks Project preamble; the block
% below is sliced MECHANICALLY, line for line, from it. Nothing was added, removed or
% redefined; the wrapper class, the xr-hyper/hyperref wiring and the externaldocument list
% are omitted because the assembler supplies the document class.
%
% lines taken from preamble.tex: 3, 5, 9, 11, 12, 156, 157, 158, ... (51 lines total)
%
% For commutative diagrams we use Xy-pic
\usepackage[all]{xy}

% We use 2cell for 2-commutative diagrams.
\xyoption{2cell}
\UseAllTwocells

\newenvironment{reference}{\comment}{\endcomment}
\newenvironment{slogan}{\comment}{\endcomment}
\newenvironment{history}{\comment}{\endcomment}
\theoremstyle{plain}
\newtheorem{theorem}[subsection]{Theorem}
\newtheorem{proposition}[subsection]{Proposition}
\newtheorem{lemma}[subsection]{Lemma}
\theoremstyle{definition}
\newtheorem{definition}[subsection]{Definition}
\newtheorem{example}[subsection]{Example}
\newtheorem{exercise}[subsection]{Exercise}
\newtheorem{situation}[subsection]{Situation}
\theoremstyle{remark}
\newtheorem{remark}[subsection]{Remark}
\newtheorem{remarks}[subsection]{Remarks}
\numberwithin{equation}{subsection}
\def\lim{\mathop{\mathrm{lim}}\nolimits}
\def\colim{\mathop{\mathrm{colim}}\nolimits}
\def\Spec{\mathop{\mathrm{Spec}}}
\def\Hom{\mathop{\mathrm{Hom}}\nolimits}
\def\Ext{\mathop{\mathrm{Ext}}\nolimits}
\def\SheafHom{\mathop{\mathcal{H}\!\mathit{om}}\nolimits}
\def\SheafExt{\mathop{\mathcal{E}\!\mathit{xt}}\nolimits}
\def\Sch{\mathit{Sch}}
\def\Mor{\mathop{\mathrm{Mor}}\nolimits}
\def\Ob{\mathop{\mathrm{Ob}}\nolimits}
\def\Sh{\mathop{\mathit{Sh}}\nolimits}
\def\NL{\mathop{N\!L}\nolimits}
\def\CH{\mathop{\mathrm{CH}}\nolimits}
\def\proetale{{pro\text{-}\acute{e}tale}}
\def\etale{{\acute{e}tale}}
\def\QCoh{\mathit{QCoh}}
\def\Ker{\mathop{\mathrm{Ker}}}
\def\Im{\mathop{\mathrm{Im}}}
\def\Coker{\mathop{\mathrm{Coker}}}
\def\Coim{\mathop{\mathrm{Coim}}}
\DeclareMathSymbol{\boxtimes}{\mathbin}{AMSa}{"02}
\def\QCohstack{\mathcal{QC}\!\mathit{oh}}
\def\Cohstack{\mathcal{C}\!\mathit{oh}}
\def\Spacesstack{\mathcal{S}\!\mathit{paces}}
\def\Quotfunctor{\mathrm{Quot}}
\def\Hilbfunctor{\mathrm{Hilb}}
\def\Curvesstack{\mathcal{C}\!\mathit{urves}}
\def\Polarizedstack{\mathcal{P}\!\mathit{olarized}}
\def\Complexesstack{\mathcal{C}\!\mathit{omplexes}}
\def\Pic{\mathop{\mathrm{Pic}}\nolimits}
\def\Picardstack{\mathcal{P}\!\mathit{ic}}
\def\Picardfunctor{\mathrm{Pic}}
\def\Deformationcategory{\mathcal{D}\!\mathit{ef}}
\usepackage{cleveref}
\crefname{originallemma}{original lemma}{original lemmas}
\Crefname{originallemma}{Original lemma}{Original lemmas}
\makeatletter
\newlabel{lemma-lisse-etale}{{0788}{source}{}{sourcepointer-1}{}}
\newlabel{lemma-lisse-etale@cref}{{[originallemma][0][]0788}{[1][0][]source}{}{}{}}
\makeatother
\crefname{originalreference}{original reference}{original references}
\Crefname{originalreference}{Original reference}{Original references}
\makeatletter
\newlabel{stacks-sheaves-lemma-cohomology-on-subcategory}{{07AK}{source}{}{sourcepointer-2}{}}
\newlabel{stacks-sheaves-lemma-cohomology-on-subcategory@cref}{{[originalreference][0][]07AK}{[1][0][]source}{}{}{}}
\makeatother
\title{Cohomology on X agrees with cohomology on the lisse-etale and flat-fppf sites}
\author{The Stacks Project Authors}
\date{stacks-project:stacks-cohomology.tex:local-snapshot-20261012; GNU FDL 1.2 or later}
\begin{document}
\maketitle
\noindent\textit{Editorial note.} All delivered mathematical text is the author's own, copied verbatim from the pinned Stacks Project chapter file stacks-cohomology.tex; no statement or proof step was rewritten, reordered or removed.\par\medskip
\noindent\textit{Allowed standard prerequisites (editorial).}\begin{itemize}
\item Cohomology on subcategories. Sheaves on Stacks, Lemma stacks-sheaves-lemma-cohomology-on-subcategory. Both parts.
\end{itemize}\medskip
\ifdefined\hypertarget\hypertarget{sourcepointer-1}{}\fi
\noindent\textit{Source reference.} Original lemma 0788: sources/stacks-chow/tags.txt, source lines 9228--9228. Cross-chapter standard prerequisite cited by the author inside the counted span; it is used as a pointer and its own proof is not part of this item.\par\medskip
\ifdefined\hypertarget\hypertarget{sourcepointer-2}{}\fi
\noindent\textit{Source reference.} Original reference 07AK: sources/stacks-chow/tags.txt, source lines 9310--9310. Cross-chapter standard prerequisite cited by the author inside the counted span; it is used as a pointer and its own proof is not part of this item.\par\medskip
\begin{lemma}
\label{lemma-lisse-etale-cohomology}
Let $\mathcal{X}$ be an algebraic stack. Notation as in
Lemma \ref{lemma-lisse-etale}.
\begin{enumerate}
\item For an abelian sheaf $\mathcal{F}$ on $\mathcal{X}_\etale$ we have
\begin{enumerate}
\item $H^p(\mathcal{X}_\etale, \mathcal{F}) =
H^p(\mathcal{X}_{lisse,\etale}, g^{-1}\mathcal{F})$, and
\item $H^p(x, \mathcal{F}) =
H^p(\mathcal{X}_{lisse,\etale}/x, g^{-1}\mathcal{F})$
for any object $x$ of $\mathcal{X}_{lisse,\etale}$.
\end{enumerate}
The same holds for sheaves of modules.
\item For an abelian sheaf $\mathcal{F}$ on $\mathcal{X}_{fppf}$ we have
\begin{enumerate}
\item $H^p(\mathcal{X}_{fppf}, \mathcal{F}) =
H^p(\mathcal{X}_{flat,fppf}, g^{-1}\mathcal{F})$, and
\item $H^p(x, \mathcal{F}) =
H^p(\mathcal{X}_{flat,fppf}/x, g^{-1}\mathcal{F})$
for any object $x$ of $\mathcal{X}_{flat,fppf}$.
\end{enumerate}
The same holds for sheaves of modules.
\end{enumerate}
\end{lemma}

\begin{proof}
Part (1)(a) follows from Sheaves on Stacks, Lemma
\ref{stacks-sheaves-lemma-cohomology-on-subcategory} applied to the inclusion
functor $\mathcal{X}_{lisse,\etale} \to \mathcal{X}_\etale$.
Part (1)(b) follows from part (1)(a). Namely, if $x$ lies over the
scheme $U$, then the site $\mathcal{X}_\etale/x$ is equivalent
to $(\Sch/U)_\etale$ and $\mathcal{X}_{lisse,\etale}$ is equivalent
to $U_{lisse,\etale}$. Part (2) is proved in the same manner.
\end{proof}

\end{document}
